\documentclass[11pt,a4paper]{article}
\usepackage[utf8]{inputenc}
\usepackage[T1]{fontenc}
\usepackage{amsmath,amsfonts,amssymb}
\usepackage{graphicx}
\usepackage{hyperref}
\usepackage{listings}
\usepackage[dvipsnames]{xcolor}
\usepackage{booktabs}
\usepackage{float}
\usepackage{geometry}
\geometry{margin=1in}

\definecolor{zenblue}{RGB}{41,121,255}
\definecolor{zengreen}{RGB}{52,199,89}
\definecolor{codegray}{RGB}{245,245,245}

\hypersetup{colorlinks=true,linkcolor=zenblue,urlcolor=zenblue,citecolor=zenblue}

\lstset{
    backgroundcolor=\color{codegray},
    basicstyle=\ttfamily\small,
    breaklines=true,
    captionpos=b,
    frame=single,
    numbers=left,
    numberstyle=\tiny\color{gray}
}

\title{
    \vspace{-2cm}
    \Large \textbf{Zen AI Model Family} \\
    \vspace{0.5cm}
    \Huge \textbf{Zen-3D} \\
    \vspace{0.3cm}
    \large Image-to-3D Generation via a Redistribution of Microsoft TRELLIS \\
    \vspace{0.5cm}
    \normalsize Technical Report v2025.01
}

\author{
    Hanzo AI Research Team\thanks{research@hanzo.ai} \and
    Zoo Labs Foundation\thanks{foundation@zoo.ngo}
}

\date{January 2025}

\begin{document}

\maketitle

\begin{abstract}
\textbf{Zen-3D} is a packaging and integration of \textbf{TRELLIS}---the open-source 3D
generation system released by Microsoft Research \cite{xiang2024trellis}---into the Zen model
stack. It is \emph{not} a from-scratch model: Zen-3D redistributes the publicly available
TRELLIS code and pretrained weights (most notably the image-conditioned
\texttt{JeffreyXiang/TRELLIS-image-large} checkpoint, $\approx$1.2B parameters) under their
original \textbf{MIT license}, wrapped behind the Zen SDK and serving interfaces. TRELLIS
generates high-quality 3D assets from a single image or a text prompt using a unified
\emph{Structured LATent} (SLAT) representation and rectified-flow transformers, and can decode
the same latent into 3D Gaussian splats, radiance fields, or textured meshes. This report
describes what TRELLIS actually is and how it works, documents which artifacts Zen-3D
redistributes and under what license, and explains the integration surface (API, output
formats, deployment) that the Zen stack adds on top. All architectural details and any reported
results are attributed to the original TRELLIS authors; we do not claim a novel architecture and
report no benchmarks of our own.
\end{abstract}

\tableofcontents
\newpage

\section{Introduction}

\subsection{What Zen-3D is, and what it is not}

Zen-3D is the Zen stack's image-to-3D and text-to-3D capability. Concretely, it is a
redistribution and integration of \textbf{TRELLIS} (``Structured 3D Latents for Scalable and
Versatile 3D Generation''), an open-source system from Microsoft Research presented at
CVPR~2025 \cite{xiang2024trellis}. The TRELLIS code and pretrained weights are released under
the \textbf{MIT license} \cite{trellis_repo}, which permits redistribution and commercial use
with attribution and license retention.

To be explicit about provenance:

\begin{itemize}
    \item Zen-3D does \textbf{not} introduce a new model architecture. There is no
    ``from-scratch'' Zen 3D network and no proprietary 3D training run behind this report.
    \item The model, weights, and generation algorithm are those of TRELLIS, authored by
    Jianfeng Xiang, Zelong Lv, Sicheng Xu, Yu Deng, Ruicheng Wang, Bowen Zhang, Dong Chen,
    Xin Tong, and Jiaolong Yang \cite{xiang2024trellis}.
    \item Zen-3D's contribution is purely at the \emph{packaging and integration} layer:
    pinning a known-good TRELLIS checkpoint, exposing it through the Zen SDK, normalizing its
    output formats, and providing serving/deployment glue.
    \item All quantitative results belong to the original work; this report reproduces no
    benchmark numbers of its own and omits any figure we cannot attribute.
\end{itemize}

\subsection{Why TRELLIS}

Single-image and text-conditioned 3D asset generation is a frontier capability for content
creation, simulation, and AR/VR. Among openly licensed options, TRELLIS is notable for (1) a
permissive MIT license on both code and weights, (2) a unified latent representation that can be
decoded into multiple downstream 3D formats (Gaussian splats, radiance fields, meshes) from a
single generation pass, and (3) strong reported quality at the 1--2B parameter scale
\cite{xiang2024trellis}. These properties make it a practical foundation to integrate rather
than a model to reinvent.

\subsection{Redistributed Artifacts}

\begin{table}[H]
\centering
\begin{tabular}{ll}
\toprule
\textbf{Property} & \textbf{Value} \\
\midrule
Upstream project & TRELLIS (Microsoft Research) \cite{xiang2024trellis,trellis_repo} \\
Task & Image-to-3D and text-to-3D generation \\
Representation & Structured LATent (SLAT) \\
Generative model & Rectified-flow transformers (two-stage) \\
Primary weights & \texttt{JeffreyXiang/TRELLIS-image-large} ($\approx$1.2B) \\
Other public weights & TRELLIS-text-base/large/xlarge (342M / 1.1B / 2.0B) \\
Output formats & 3D Gaussian splats, radiance fields, textured meshes \\
License (code + weights) & MIT \cite{trellis_repo} \\
Training data (upstream) & TRELLIS-500K ($\approx$500K curated 3D assets) \\
Zen-3D additions & SDK wrapper, serving, output normalization \\
\bottomrule
\end{tabular}
\caption{Zen-3D: artifacts redistributed from TRELLIS and the integration layer added by Zen.
Parameter counts and dataset sizes are as reported by the TRELLIS authors
\cite{xiang2024trellis}.}
\end{table}

\section{TRELLIS Architecture}

This section summarizes the TRELLIS method as described by its authors \cite{xiang2024trellis}.
It is included for completeness; none of it is original to Zen-3D.

\subsection{Structured LATent (SLAT) Representation}

The cornerstone of TRELLIS is the \emph{Structured LATent} (SLAT) representation: a set of local
latent vectors attached to the active (occupied) voxels of a sparse 3D grid. Formally, an asset
is encoded as
\begin{equation}
    \mathbf{z} = \{(\boldsymbol{z}_i, \boldsymbol{p}_i)\}_{i=1}^{L},
    \qquad \boldsymbol{p}_i \in \{0,\dots,N-1\}^3,
\end{equation}
where each $\boldsymbol{p}_i$ is the integer coordinate of an active voxel and $\boldsymbol{z}_i$
is its local latent vector. The active voxels $\boldsymbol{p}_i$ outline the coarse geometry of
the asset, while the latents $\boldsymbol{z}_i$ capture finer structural and textural detail. By
default the grid resolution is $N = 64$, which yields roughly $L \approx 20{,}000$ active voxels
per asset on average \cite{xiang2024trellis}.

The per-voxel features used to fit the latents are obtained from 2D vision features rather than
from raw geometry. TRELLIS renders an asset from many camera views sampled on a sphere, extracts
feature maps with a pretrained \textbf{DINOv2} encoder, projects each active voxel into those
multiview feature maps, and averages the retrieved features to form the voxel feature
$\boldsymbol{f}_i$ \cite{xiang2024trellis}. This couples geometry (the sparse grid) with
appearance (the aggregated visual features) in a single structured latent.

\subsection{SLAT Variational Autoencoder and Decoders}

A transformer-based VAE encodes the voxel features into SLAT and reconstructs 3D outputs from
it. Voxel tokens are serialized with sinusoidal positional encodings derived from their
coordinates. A distinctive property of TRELLIS is that the \emph{same} SLAT can be decoded into
several 3D output formats by separate, format-specific decoders \cite{xiang2024trellis}:

\begin{itemize}
    \item \textbf{3D Gaussians.} Each latent decodes to a small set of Gaussians with position
    offsets, colors, scales, opacities, and rotations. Positions are anchored to the voxel and
    perturbed by a bounded offset, e.g.\ $\boldsymbol{x}_i^k = \boldsymbol{p}_i +
    \tanh(\boldsymbol{o}_i^k)$. The encoder and the Gaussian decoder are trained end-to-end.
    \item \textbf{Radiance fields.} Each latent decodes to a local radiance volume represented
    with a CP-decomposition (following Strivec), enabling neural volumetric rendering.
    \item \textbf{Meshes.} Each latent decodes to FlexiCubes parameters and signed-distance
    values, from which a textured surface mesh is extracted (upsampled to a higher resolution).
\end{itemize}

The non-Gaussian decoders are trained separately on top of a frozen encoder
\cite{xiang2024trellis}.

\subsection{Two-Stage Rectified-Flow Generation}

TRELLIS generates SLAT with two \textbf{rectified-flow transformers} trained with a conditional
flow-matching objective, run in sequence \cite{xiang2024trellis}:

\begin{enumerate}
    \item \textbf{Sparse-structure stage.} A flow transformer $\mathcal{G}_S$ generates the
    sparse occupancy structure---i.e.\ which voxels are active. The binary grid is first
    compressed by a VAE into a low-resolution dense feature grid, and the flow model operates in
    that space before decoding back to the occupied voxel set.
    \item \textbf{Structured-latent stage.} A second flow transformer $\mathcal{G}_L$ generates
    the local latents $\boldsymbol{z}_i$ for the active voxels produced by stage~1, operating
    over the sparse set of occupied cells.
\end{enumerate}

Both transformers use attention adapted to 3D sparsity (e.g.\ serialization of active voxels and
shifted-window attention in 3D); $\mathcal{G}_S$ and $\mathcal{G}_L$ are trained separately
\cite{xiang2024trellis}.

\subsection{Conditioning}

TRELLIS supports both image and text conditioning, injected into the flow transformers via
cross-attention \cite{xiang2024trellis}:

\begin{itemize}
    \item \textbf{Image conditioning} uses DINOv2 visual features of the input image. This is
    the path exercised by the \texttt{TRELLIS-image-large} checkpoint that Zen-3D primarily
    redistributes.
    \item \textbf{Text conditioning} uses CLIP text features, corresponding to the
    text-conditioned TRELLIS checkpoints.
\end{itemize}

\subsection{Model Sizes and Training Data (as reported upstream)}

The TRELLIS authors report training models at three scales---\textbf{342M} (Base), \textbf{1.1B}
(Large), and \textbf{2.0B} (X-Large)---on \textbf{TRELLIS-500K}, a curated set of roughly 500K
3D assets sourced from Objaverse(XL), ABO, 3D-FUTURE, HSSD, and Toys4k and filtered by aesthetic
score \cite{xiang2024trellis,trellis_repo}. The image-conditioned checkpoint redistributed by
Zen-3D, \texttt{JeffreyXiang/TRELLIS-image-large}, is the $\approx$1.2B-parameter image variant.

\begin{table}[H]
\centering
\begin{tabular}{lll}
\toprule
\textbf{Public checkpoint} & \textbf{Parameters} & \textbf{Conditioning} \\
\midrule
TRELLIS-image-large & $\approx$1.2B & Image \\
TRELLIS-text-base & 342M & Text \\
TRELLIS-text-large & 1.1B & Text \\
TRELLIS-text-xlarge & 2.0B & Text \\
\bottomrule
\end{tabular}
\caption{Publicly released TRELLIS checkpoints, as listed by the upstream project
\cite{trellis_repo}. Zen-3D primarily integrates the image-conditioned model.}
\end{table}

\section{Reported Results (Upstream)}

We do not run our own evaluation, and we do not restate numerical benchmarks here because we
cannot independently verify them in this report. The TRELLIS paper reports that the method
produces diverse, high-quality 3D assets and compares favorably against prior and contemporaneous
3D generators at similar scale; for the exact quantitative comparisons, evaluation protocol, and
qualitative galleries, we refer the reader to the original publication and project page
\cite{xiang2024trellis,trellis_repo}. Any number not attributable to that source should be
considered absent rather than implied.

\section{Zen Integration}

Everything in this section is the integration layer that Zen-3D adds; it does not modify the
TRELLIS model itself.

\subsection{What Zen-3D Provides}

\begin{itemize}
    \item \textbf{Pinned redistribution.} A fixed, known-good TRELLIS checkpoint and code
    revision, repackaged under the Zen distribution while retaining the upstream MIT license and
    attribution (see Section~\ref{sec:licensing}).
    \item \textbf{SDK surface.} A thin Python wrapper that loads the upstream
    \texttt{TrellisImageTo3DPipeline} and exposes a uniform Zen API for image-to-3D (and
    text-to-3D where the corresponding checkpoint is used).
    \item \textbf{Output normalization.} Helpers to export the decoded SLAT into the format a
    caller wants---3D Gaussian splats, a radiance field, or a textured mesh (e.g.\ GLB/OBJ)---using
    TRELLIS's own decoders.
    \item \textbf{Serving glue.} Deployment configuration for running the pipeline behind the Zen
    serving stack.
\end{itemize}

\subsection{Output Formats}

Because a single SLAT can be decoded multiple ways, the same generation can yield different
deliverables without re-running the flow models:

\begin{table}[H]
\centering
\begin{tabular}{lll}
\toprule
\textbf{Output} & \textbf{Decoder (TRELLIS)} & \textbf{Typical use} \\
\midrule
3D Gaussian splats & Gaussian decoder & Fast novel-view rendering \\
Radiance field & CP-decomposition volume & Volumetric rendering \\
Textured mesh & FlexiCubes + SDF & DCC tools, game/CAD pipelines \\
\bottomrule
\end{tabular}
\caption{TRELLIS output formats exposed through the Zen-3D API. Decoders are those of the
upstream model \cite{xiang2024trellis}.}
\end{table}

\subsection{Python API}

The Zen-3D wrapper mirrors the upstream TRELLIS pipeline. Usage is illustrative:

\begin{lstlisting}[language=Python, caption=Zen-3D image-to-3D (wrapping TRELLIS)]
from zen import Zen3D
from PIL import Image

# Loads the redistributed TRELLIS-image-large checkpoint (MIT licensed).
model = Zen3D.from_pretrained("zenlm/zen-3d")

image = Image.open("input.png")
outputs = model.generate(image)

# Decode the same SLAT into the format you need.
outputs.gaussians.save("asset.ply")    # 3D Gaussian splats
outputs.mesh.export("asset.glb")       # textured mesh (FlexiCubes)
# outputs.radiance_field for volumetric rendering
\end{lstlisting}

\subsection{Deployment Notes}

TRELLIS depends on differentiable rendering submodules (e.g.\ a differentiable octree renderer
and a modified FlexiCubes) for training and for certain decoding paths; these submodules carry
their own licenses (see Section~\ref{sec:licensing}) and have their own build requirements.
Inference for the redistributed image model runs on a single modern GPU; for exact hardware
requirements, dependency versions, and build steps, follow the upstream project documentation
\cite{trellis_repo}.

\section{Licensing and Attribution}
\label{sec:licensing}

\begin{itemize}
    \item \textbf{TRELLIS code and weights: MIT.} The TRELLIS repository and the released
    checkpoints (including \texttt{TRELLIS-image-large}) are MIT licensed
    \cite{trellis_repo}. Zen-3D's redistribution retains the upstream MIT license text and
    copyright notice, and credits Microsoft Research and the TRELLIS authors.
    \item \textbf{Submodule exceptions.} Some optional submodules are not MIT: the
    differentiable octree renderer (\texttt{diffoctreerast}) and the modified FlexiCubes are
    licensed under their own terms \cite{trellis_repo}. Deployments that build these components
    must comply with those licenses.
    \item \textbf{Upstream dependencies.} The vision/text encoders used for conditioning
    (DINOv2, CLIP) and the data sources behind TRELLIS-500K carry their own licenses and
    usage terms, which downstream users should review independently.
    \item \textbf{No relicensing of capabilities.} Zen-3D claims no ownership of the TRELLIS
    model or its outputs' underlying method. Where this report describes architecture or
    results, the source of record is the TRELLIS paper and repository
    \cite{xiang2024trellis,trellis_repo}.
\end{itemize}

\section{Conclusion}

Zen-3D is an honest packaging of Microsoft's TRELLIS into the Zen stack: an MIT-licensed,
image-conditioned (and optionally text-conditioned) 3D generator built on the Structured LATent
representation and two-stage rectified-flow transformers, capable of decoding to 3D Gaussian
splats, radiance fields, or meshes. The model, its architecture, and any reported quality are
the work of the TRELLIS authors \cite{xiang2024trellis}; Zen-3D adds only redistribution, an SDK
wrapper, output normalization, and serving integration. We make no from-scratch architecture or
benchmark claims, and we encourage users to consult the upstream project for authoritative
technical detail and evaluation.

\begin{thebibliography}{10}
\bibitem{xiang2024trellis} J. Xiang, Z. Lv, S. Xu, Y. Deng, R. Wang, B. Zhang, D. Chen, X. Tong, and J. Yang, ``Structured 3D Latents for Scalable and Versatile 3D Generation,'' CVPR, 2025 (arXiv:2412.01506). \url{https://arxiv.org/abs/2412.01506}
\bibitem{trellis_repo} Microsoft Research, ``TRELLIS,'' GitHub repository (MIT license; weights \texttt{JeffreyXiang/TRELLIS-image-large}). \url{https://github.com/microsoft/TRELLIS}
\end{thebibliography}

\end{document}
